\section*{2026-06-01}

\subsection*{Python 贪吃蛇}

pip install pygame \par

\lstset{language=bash}
\begin{lstlisting}
snake_game/
--- main.py          # 程序入口
--- config.py        # 全局配置(尺寸、颜色、速度)
--- snake.py         # 蛇类封装
--- food.py          # 食物类封装
--- game.py          # 游戏逻辑主类
\end{lstlisting}

\subsubsection*{config.py}
\lstset{language=Python}
\begin{lstlisting}
# 窗口与格子配置
WIDTH = 600
HEIGHT = 400
BLOCK_SIZE = 20
SPEED = 12

# 颜色 RGB
BLACK = (0, 0, 0)
WHITE = (255, 255, 255)
RED = (255, 0, 0)
GREEN = (0, 255, 0)
BLUE = (0, 0, 255)
\end{lstlisting}

\subsubsection*{snake.py}
\lstset{language=Python}
\begin{lstlisting}
import pygame
from config import BLOCK_SIZE, GREEN

class Snake:
    def __init__(self):
        self.body = [(100,100), (80,100), (60,100)]  # 身体坐标列表
        self.dir = (BLOCK_SIZE, 0)  # 初始向右

    def change_dir(self, new_dir):
        # 禁止反向掉头
        if (new_dir[0] != -self.dir[0]) or (new_dir[1] != -self.dir[1]):
            self.dir = new_dir

    def move(self, grow=False):
        head_x, head_y = self.body[0]
        new_head = (head_x + self.dir[0], head_y + self.dir[1])
        self.body.insert(0, new_head)
        if not grow:
            self.body.pop()

    def draw(self, screen):
        for pos in self.body:
            rect = pygame.Rect(pos[0], pos[1], BLOCK_SIZE-1, BLOCK_SIZE-1)
            pygame.draw.rect(screen, GREEN, rect)

    def get_head(self):
        return self.body[0]
\end{lstlisting}

\subsubsection*{food.py}
\lstset{language=Python}
\begin{lstlisting}
import random
import pygame
from config import BLOCK_SIZE, WIDTH, HEIGHT, RED

class Food:
    def __init__(self):
        self.reset_pos()

    def reset_pos(self):
        # 随机生成在格子上
        self.x = random.randrange(0, WIDTH-BLOCK_SIZE, BLOCK_SIZE)
        self.y = random.randrange(0, HEIGHT-BLOCK_SIZE, BLOCK_SIZE)
        self.pos = (self.x, self.y)

    def draw(self, screen):
        rect = pygame.Rect(self.x, self.y, BLOCK_SIZE-1, BLOCK_SIZE-1)
        pygame.draw.rect(screen, RED, rect)
\end{lstlisting}

\subsubsection*{game.py}
\lstset{language=Python}
\begin{lstlisting}
import pygame
from snake import Snake
from food import Food
from config import *

class Game:
    def __init__(self):
        pygame.init()
        self.screen = pygame.display.set_mode((WIDTH, HEIGHT))
        pygame.display.set_caption("经典贪吃蛇")
        self.clock = pygame.time.Clock()
        self.font = pygame.font.SysFont(None, 35)
        self.snake = Snake()
        self.food = Food()
        self.score = 0
        self.game_over = False

    def check_collision(self):
        head = self.snake.get_head()
        # 撞墙
        if head[0]<0 or head[0]>=WIDTH or head[1]<0 or head[1]>=HEIGHT:
            return True
        # 撞到自己身体
        if head in self.snake.body[1:]:
            return True
        return False

    def eat_food(self):
        if self.snake.get_head() == self.food.pos:
            self.score +=10
            self.food.reset_pos()
            return True
        return False

    def event_handle(self):
        for event in pygame.event.get():
            if event.type == pygame.QUIT:
                pygame.quit()
                exit()
            if event.type == pygame.KEYDOWN:
                if event.key == pygame.K_UP:
                    self.snake.change_dir((0, -BLOCK_SIZE))
                elif event.key == pygame.K_DOWN:
                    self.snake.change_dir((0, BLOCK_SIZE))
                elif event.key == pygame.K_LEFT:
                    self.snake.change_dir((-BLOCK_SIZE, 0))
                elif event.key == pygame.K_RIGHT:
                    self.snake.change_dir((BLOCK_SIZE, 0))

    def draw_ui(self):
        text = self.font.render(f"Score:{self.score}", True, WHITE)
        self.screen.blit(text, [0,0])

    def run(self):
        while not self.game_over:
            self.event_handle()
            grow = self.eat_food()
            self.snake.move(grow)

            if self.check_collision():
                self.game_over = True

            self.screen.fill(BLACK)
            self.snake.draw(self.screen)
            self.food.draw(self.screen)
            self.draw_ui()

            pygame.display.update()
            self.clock.tick(SPEED)

        # 游戏结束弹窗
        over_text = self.font.render("Game Over! 关闭窗口退出", True, WHITE)
        self.screen.blit(over_text, [WIDTH//6, HEIGHT//2])
        pygame.display.update()
        while True:
            for e in pygame.event.get():
                if e.type == pygame.QUIT:
                    pygame.quit()
                    return
\end{lstlisting}

\subsubsection*{main.py}
\lstset{language=Python}
\begin{lstlisting}
from game import Game

if __name__ == "__main__":
    game = Game()
    game.run()

\end{lstlisting}

